% Include after the Clifford-rotation theorem; install only after replay passes.
\paragraph{The additional multiplicity-four exclusion $R6\_4\_174$.}
The four invertible basis elements have Frobenius--Perron dimension $1$;
the remaining two have dimension $d=4+3\sqrt2$, with $d^2=2+8d$.
Thus the global dimension is $72+48\sqrt2$, whose norm is $576$.
The involution-module calculation has a single antiinvariant branch, and
its stabilizer commutator is negative on both noninvertible simples.
Each cubic invariant space has dimension four. Since the global-dimension
prime support contains $3$, the exact relative traces of their third
indicators are allowed to be $-2,-1/2,$ or $1$. The computation first uses
the weaker interval $[-2,1]$.

At trace height six, there are ten low induction frames. Eight character
labellings survive the earlier exact conditions; their full second-indicator
assignments give 128 systems. Every system admits an exact separating ray,
already with the interval constraint. Hence no pseudo-unitary
categorification exists. The two distinct Frobenius--Perron dimensions
supply the previously established general-to-Galois-pseudo-unitary reduction,
so the exclusion holds for arbitrary complex fusion categories.
The composite certificate uses the general higher-stabilizer verifier,
with the reduction payload \texttt{two\_fp\_values}; all matrix and branch
data are reconstructed from the input tensor.
